\documentclass[12pt]{article}
\usepackage[margin=0.75in]{geometry}
\usepackage{fontspec}
\setmainfont{Latin Modern Roman}
\usepackage{microtype}
\usepackage{fancyhdr}
\usepackage{titlesec}
\usepackage{enumitem}
\usepackage{xcolor}
\usepackage[hidelinks]{hyperref}
\setlength{\parindent}{1.1em}
\setlength{\parskip}{0.38em}
\setlength{\headheight}{15pt}
\titleformat{\section}{\large\bfseries\scshape}{}{0pt}{}
\titleformat{\subsection}{\normalsize\bfseries}{}{0pt}{}
\pagestyle{fancy}
\fancyhf{}
\fancyfoot[C]{\thepage}
\renewcommand{\headrulewidth}{0.4pt}
\newcommand{\focusbox}[1]{\par\vspace{0.4em}\noindent\begingroup\setlength{\fboxsep}{6pt}\colorbox{gray!12}{\begin{minipage}{\dimexpr\linewidth-2\fboxsep\relax}#1\end{minipage}}\endgroup\par\vspace{0.5em}}
\newcommand{\studentline}{\vspace{0.35em}\hrule\vspace{0.65em}}
\newcommand{\shortline}{\vspace{0.25em}\hrule\vspace{0.45em}}
\newcommand{\longspace}{\vspace{0.95in}}

\fancyhead[L]{Whately: Syllogism Lab}
\fancyhead[R]{Student Handout}
\begin{document}
\begin{center}
{\LARGE\bfseries Week 4 Argument Lab}\par\vspace{0.25em}
{\large\scshape Practical Syllogism Maps}\par\vspace{0.6em}
\rule{0.82\textwidth}{0.4pt}
\end{center}

\section*{Name: \rule{2.4in}{0.4pt} \hfill Date: \rule{1.3in}{0.4pt}}

\section*{Purpose}
This lab turns Week 4 reading into a portfolio artifact. You will map practical syllogisms by finding the premises, conclusion, major term, minor term, and middle term. The goal is not yet full mood-and-figure classification. The goal is to see whether the middle term really connects the conclusion.

\section*{Part I: Model}
\textbf{Argument:}
\begin{itemize}[leftmargin=*]
\item All students who submit late work without an approved extension should receive a penalty.
\item Marcus submitted late work without an approved extension.
\item Therefore, Marcus should receive a penalty.
\end{itemize}

\begin{itemize}[leftmargin=*]
\item \textbf{Conclusion:} Marcus should receive a penalty.
\item \textbf{Minor term:} Marcus.
\item \textbf{Major term:} people who should receive a penalty.
\item \textbf{Middle term:} students who submit late work without an approved extension.
\item \textbf{Major premise:} All students who submit late work without an approved extension should receive a penalty.
\item \textbf{Minor premise:} Marcus submitted late work without an approved extension.
\item \textbf{Does the middle term connect?} Yes. If both premises are true, the conclusion follows.
\end{itemize}

\section*{Part II: Map the Syllogism}
For each argument, identify the conclusion, major term, minor term, middle term, major premise, and minor premise. Then decide whether the conclusion follows from the premises.

\subsection*{1. Strong Form}
\textbf{Argument:} All unjust laws deserve public resistance. This law is unjust. Therefore, this law deserves public resistance.

\textbf{Conclusion:}\shortline
\textbf{Major term:}\shortline
\textbf{Minor term:}\shortline
\textbf{Middle term:}\shortline
\textbf{Major premise:}\shortline
\textbf{Minor premise:}\shortline
\textbf{Does the conclusion follow if the premises are true? Explain.}\studentline

\subsection*{2. Tempting but Weak}
\textbf{Argument:} All careful writers revise their sentences. Elena revises her sentences. Therefore, Elena is a careful writer.

\textbf{Conclusion:}\shortline
\textbf{Major term:}\shortline
\textbf{Minor term:}\shortline
\textbf{Middle term:}\shortline
\textbf{Major premise:}\shortline
\textbf{Minor premise:}\shortline
\textbf{Does the conclusion follow if the premises are true? Explain.}\studentline

\subsection*{3. Hidden Premise}
\textbf{Argument as spoken:} "This speech makes the crowd angry, so it must be politically dangerous."

\textbf{Supply the missing major premise:}\studentline
\textbf{Minor premise:}\shortline
\textbf{Conclusion:}\shortline
\textbf{Middle term:}\shortline
\textbf{Does the supplied premise make the argument stronger, weaker, or more questionable?}\studentline

\section*{Part III: Shakespeare Application}
Choose \textbf{one} of the two literary cases below and create a practical syllogism map. Mark this section with a star if you want it considered for your Logic Portfolio.

\subsection*{Option A: Brutus in \textit{Julius Caesar}}
Brutus reasons that Caesar may become dangerous if crowned, and that Rome must act before the danger fully appears.

\textbf{Possible conclusion:} Caesar should be stopped before he is crowned.

\textbf{Major premise:}\studentline
\textbf{Minor premise:}\studentline
\textbf{Conclusion:}\shortline
\textbf{Major term:}\shortline
\textbf{Minor term:}\shortline
\textbf{Middle term:}\shortline
\textbf{Which premise is most debatable, and why?}\studentline\studentline

\subsection*{Option B: Macbeth and Prophecy}
Macbeth hears the witches' prediction and begins to think about the crown as if prophecy changes the meaning of his choices.

\textbf{Possible conclusion:} Macbeth's path to the crown is necessary rather than blameworthy.

\textbf{Major premise:}\studentline
\textbf{Minor premise:}\studentline
\textbf{Conclusion:}\shortline
\textbf{Major term:}\shortline
\textbf{Minor term:}\shortline
\textbf{Middle term:}\shortline
\textbf{Which premise is most debatable, and why?}\studentline\studentline

\section*{Part IV: Repair a Faulty Syllogism}
Choose one faulty or questionable argument from this lab. Rewrite it so that either (1) the conclusion follows more carefully, or (2) the conclusion is narrowed to fit the premises.

\textbf{Original argument I am repairing:}\studentline
\textbf{Problem with the original middle term or premise:}\studentline
\textbf{Repaired version:}\studentline\studentline

\section*{Part V: Window Note and Portfolio Summary}
Create a four-window note titled \textbf{Two Claims and a Conclusion}. Use drawings, symbols, arrows, or a simple diagram. Your four windows should show:
\begin{enumerate}[leftmargin=*]
\item the major premise;
\item the minor premise;
\item the middle term as a bridge;
\item the conclusion, with a check mark if it follows or a warning sign if it does not.
\end{enumerate}

At the bottom of the window note, write a 6--8 sentence portfolio summary explaining how a syllogism map helps a reader test an argument instead of merely reacting to its persuasive force. Use one example from the lab, \textit{Macbeth}, or \textit{Julius Caesar}.

\focusbox{\textbf{Portfolio summary starter:} A syllogism map helps me test an argument because \rule{2.2in}{0.4pt}. The most important term to watch is the middle term because \rule{1.8in}{0.4pt}.}

\vspace{1.1in}
\end{document}
